\section{Minimum norm of regularized regional filters}
\label{sec:peps_regularized_patch_minimum}

Let $V$ be a finite ordered set of sites, with finite-dimensional coordinate
Hilbert spaces $\mathcal H_v$, and set
$\mathcal H_X=\bigotimes_{v\in X}\mathcal H_v$.
The empty tensor product is $\C$.
For $X\subseteq V$, use the regional extension $\widehat K_X$ of
Definition~\ref{def:peps_regional_operator_extension}.
Fix regions $X_0,\ldots,X_{m-1}$ and a unit vector $\Omega\in\mathcal H_V$.
The following argument proves attainment and the elementary norm bounds in the
\href{https://github.com/openai/math/blob/adc7f1241b42e322a6451854ab7e4b4c146bf78a/preprints/Polynomial-PEPS-approximation-of-gapped-square-grid-ground-states-September-24-2026/build/sections/03-patches.tex\#L68-L99}{regularized variational problem}
of the September 24, 2026 polynomial-PEPS manuscript.
% Source: 03-patches.tex, lines 51--114, commit
% adc7f1241b42e322a6451854ab7e4b4c146bf78a.
% These are the finite-dimensional minimum and elementary bounds only, not
% stationarity, energy, slow regulator growth, retained rank, or Proposition 4.1.

\begin{definition}[Feasible densities and ordered filters]
    \label{def:peps_regularized_patch_objective}
    \lean{
        TNLean.PEPS.regularizedPatchDomain,
        TNLean.PEPS.regularizedPatchFilter,
        TNLean.PEPS.regularizedPatchOperator,
        TNLean.PEPS.regularizedPatchOutput,
        TNLean.PEPS.regularizedPatchObjective
    }
    \leanok
    \uses{def:peps_regional_operator_extension}
    Let $b>0$ and $a_0,\ldots,a_{m-1}\in\R$, and put $A=\sum_j a_j$. Define
    \begin{align*}
        \mathcal D & =\prod_{j=0}^{m-1}
        \{x_j\in\operatorname{End}(\mathcal H_{X_j}):x_j\ge0,\ \tr x_j=1\}, \\
        K_j(x)     & =(x_j+bI_{X_j})^{-a_j/2},\qquad L_j(x)=\widehat{K_j(x)}_{X_j}, \\
        M(x)       & =L_{m-1}(x)\cdots L_0(x),                              \\
        \Psi(x)    & =M(x)\Omega,\qquad F(x)=\|\Psi(x)\|.
    \end{align*}
    The densities are independent for different indices, including indices
    with the same region. The powers act on the full regional space.
    For $m=0$, the product is the identity and $F=1$.

    In the special case of two occurrences $X_0=X_1=X$, write
    $K_j=(x_j+bI_X)^{-a_j/2}$. The upper line below carries $\mathcal H_X$;
    the lower line carries the complementary tensor factor
    $\mathcal H_{V\setminus X}$, which may be one-dimensional.
    This picture illustrates repeated support only; the general family of
    regions remains arbitrary.
    % Formula: ((K_1 tensor I)(K_0 tensor I) Omega)_{i,alpha}
    % = sum_{j,k} (K_1)_{i,j}(K_0)_{j,k} Omega_{k,alpha}.
    % Source: 03-patches.tex lines 73--90 at adc7f1241b42e322a6451854ab7e4b4c146bf78a.
    % Ink-to-index map: upper row is the repeated regional physical space;
    % lower row is its complement. Omega has those two physical output legs.
    % Each K_j acts only on the upper row; the aligned lower I is the
    % complementary identity in that factor. Action starts at Omega on the right.
    % Contractions: the four internal physical wires; no swaps or crossings.
    % Open legs: i on the upper left and alpha on the lower left.
    % Boundary signature: phys:w, phys:w; output lies in H_X tensor H_(V\X).
    % This repeated-region example assumes no commutation and asserts no nesting
    % condition for the general family of regions.
    \begin{center}
        \tenkzkernel
        \begin{tenkz}[rows={wire,wire}, cols=5, bonds=none]
            \tn[at=(1,1), name=filterone, skin=box,
                ports={180:physical, 0:physical}]{K_1}
            \tn[at=(1,3), name=filterzero, skin=box,
                ports={180:physical, 0:physical}]{K_0}
            \tn[at=(2,1), name=identityone, skin=box,
                ports={180:physical, 0:physical}]{I}
            \tn[at=(2,3), name=identityzero, skin=box,
                ports={180:physical, 0:physical}]{I}
            \tn[at=(1,5), name=reference, skin=box, wires=2,
                ports={180@1:physical, 180@2:physical}]{\Omega}
            \tnwire[name=insideoutput]{open w}{filterone.180}
            \tnwire{filterone.0}{filterzero.180}
            \tnwire{filterzero.0}{reference.180@1}
            \tnwire[name=outsideoutput]{open w}{identityone.180}
            \tnwire{identityone.0}{identityzero.180}
            \tnwire{identityzero.0}{reference.180@2}
            \tnmark[form=label, label pos=n]{on insideoutput 0.5}{$i$}
            \tnmark[form=label, label pos=s]{on outsideoutput 0.5}{$\alpha$}
        \end{tenkz}
    \end{center}
    Reading from $\Omega$ on the right, $K_0\otimes I$ acts first and
    $K_1\otimes I$ acts second. Thus the displayed output is
    $\Psi(x)=(K_1\otimes I)(K_0\otimes I)\Omega$.
\end{definition}

For a qubit example, take $X=V$, $b=1$, $a_0=a_1=2$,
$x_0=\ketbra{0}{0}$ and $x_1=\ketbra{+}{+}$, where
$\ket{+}=(\ket{0}+\ket{1})/\sqrt2$. Then
\[
    K_0=\begin{pmatrix}1/2&0\\0&1\end{pmatrix},\qquad
    K_1=\begin{pmatrix}3/4&-1/4\\-1/4&3/4\end{pmatrix}.
\]
The product order matters:
$(K_1K_0)_{01}=-1/4$, whereas $(K_0K_1)_{01}=-1/8$.
The complementary space is $\C$ in this example, so the lower identities
are scalar identities.


\begin{lemma}[Compactness and continuity]
    \label{lem:peps_regularized_patch_compact_continuous}
    \lean{
        TNLean.PEPS.isCompact_regularizedPatchDomain,
        TNLean.PEPS.regularizedPatchDomain_nonempty,
        TNLean.PEPS.continuousOn_regularizedPatchFilter,
        TNLean.PEPS.continuousOn_regularizedPatchObjective
    }
    \leanok
    \uses{def:peps_regularized_patch_objective}
    The set $\mathcal D$ is compact and nonempty. Each $L_j$ and the
    objective $F$ are continuous on $\mathcal D$.
\end{lemma}
\begin{proof}
    \leanok
    \uses{thm:peps_regional_extension_norm}
    In finite dimension, the positive semidefinite trace-one matrices form a
    compact set. Their finite product is compact. Since $\|\Omega\|=1$,
    some global configuration has a nonzero coefficient in $\Omega$.
    Restricting that configuration to each region gives
    $d_j=\dim\mathcal H_{X_j}>0$. The densities $I_{X_j}/d_j$ prove
    nonemptiness. For $b>0$, the shifted matrices are positive definite,
    and their fixed real powers vary continuously. Identity extension,
    finite ordered multiplication, evaluation on $\Omega$, and the norm
    preserve continuity.
\end{proof}

\begin{theorem}[Bounds for every feasible tuple]
    \label{thm:peps_regularized_patch_bounds}
    \lean{
        TNLean.PEPS.regularizedPatchFilter_norm_bounds,
        TNLean.PEPS.regularizedPatchObjective_bounds,
        TNLean.PEPS.regularizedPatchObjective_pos
    }
    \leanok
    \uses{def:peps_regularized_patch_objective}
    Suppose $a_j\ge0$ for every $j$.
    For every $x\in\mathcal D$, every $j$, and every $\xi\in\mathcal H_V$,
    \[
        (1+b)^{-a_j/2}\|\xi\|
        \le\|L_j(x)\xi\|\le b^{-a_j/2}\|\xi\|.
    \]
    Consequently, every $x\in\mathcal D$ satisfies
    \[
        0<(1+b)^{-A/2}\le F(x)\le b^{-A/2}.
    \]
\end{theorem}
\begin{proof}
    \leanok
    \uses{thm:peps_regional_extension_norm}
    A positive semidefinite trace-one matrix satisfies $0\le x_j\le I$.
    The eigenvalues of $x_j+bI$ therefore lie in $[b,1+b]$.
    The negative power has operator norm at most $b^{-a_j/2}$, while its
    inverse has norm at most $(1+b)^{a_j/2}$.
    Apply the identity-extension inequality to the factor and its inverse
    to obtain the upper and lower bounds. Apply these bounds successively
    in the order $L_0,\ldots,L_{m-1}$, and use $\|\Omega\|=1$.
    This argument never exchanges two factors.
\end{proof}

\begin{definition}[Minimum and normalization]
    \label{def:peps_regularized_patch_minimum}
    \lean{
        TNLean.PEPS.regularizedPatchMinimum,
        TNLean.PEPS.normalizedRegularizedPatchOutput
    }
    \leanok
    \uses{def:peps_regularized_patch_objective}
    Define the minimum value and normalized output by
    \[
        N=\inf_{x\in\mathcal D}F(x),\qquad
        \widehat\Psi(x)=F(x)^{-1}\Psi(x).
    \]
    Under the nonnegative-weight hypothesis, the preceding theorem makes
    this normalization well defined for every feasible tuple.
\end{definition}

\begin{theorem}[Attainment and common value]
    \label{thm:peps_regularized_patch_minimum}
    \lean{
        TNLean.PEPS.exists_isMinOn_regularizedPatchObjective,
        TNLean.PEPS.regularizedPatchObjective_eq_minimum,
        TNLean.PEPS.regularizedPatchMinimum_bounds,
        TNLean.PEPS.regularizedPatchMinimum_pos,
        TNLean.PEPS.norm_normalizedRegularizedPatchOutput,
        TNLean.PEPS.normalizedRegularizedPatchOutput_eq_minimum_smul
    }
    \leanok
    \uses{def:peps_regularized_patch_minimum}
    There exists $x_*\in\mathcal D$ such that $F(x_*)\le F(x)$ for all
    $x\in\mathcal D$. Every feasible minimizer has objective value $N$.
    If all $a_j\ge0$, then
    \[
        0<(1+b)^{-A/2}\le N\le b^{-A/2},\qquad
        \|\widehat\Psi(x)\|=1\quad(x\in\mathcal D).
    \]
    At every feasible minimizer $x_*$,
    $\widehat\Psi(x_*)=N^{-1}M(x_*)\Omega$.
\end{theorem}
\begin{proof}
    \leanok
    \uses{lem:peps_regularized_patch_compact_continuous,thm:peps_regularized_patch_bounds}
    A continuous real-valued function on a nonempty compact set attains
    its minimum. The value at any minimizing point is the least element
    of $F(\mathcal D)$, hence its infimum. Apply the bounds for every
    feasible tuple at a minimizer. Positivity permits division by $F(x)$;
    at a minimizer this divisor equals $N$.
\end{proof}

\begin{corollary}[Exponential regularization]
    \label{cor:peps_regularized_patch_exp_bounds}
    \lean{
        TNLean.PEPS.regularizedPatchObjective_exp_bounds,
        TNLean.PEPS.regularizedPatchMinimum_exp_bounds
    }
    \leanok
    \uses{def:peps_regularized_patch_minimum}
    Assume all $a_j\ge0$ and set $b=e^{-R}$. For every real $R$ and
    every feasible tuple $x$,
    \[
        (1+e^{-R})^{-A/2}\le F(x)\le e^{RA/2}.
    \]
    For $R\ge0$, the minimum satisfies
    \[
        2^{-A/2}\le N\le e^{RA/2}.
    \]
\end{corollary}
\begin{proof}
    \leanok
    \uses{thm:peps_regularized_patch_bounds,thm:peps_regularized_patch_minimum}
    Substitute $b=e^{-R}>0$. If $R\ge0$, then $1+e^{-R}\le2$,
    so the nonpositive exponent $-A/2$ gives
    $(1+e^{-R})^{-A/2}\ge2^{-A/2}$.
\end{proof}

The argument allows empty or repeated regions, zero weights, and singular
densities. Neither nestedness of the regions nor commutation of their
filters is required. It gives the same minimum value for every minimizer,
without asserting uniqueness or smooth dependence of a minimizer on $R$.
